\section{Causal simulation, calibration and the same task}\label{sec:composition}
\begin{proposition}[Typed resource and risk composition]\label{prop:morphism}
Fix the protected audit protocol, its response space and a loss in $[0,L]$. Suppose a parameter-independent causal simulator from a source experiment to a target has $R$ internal labels, legal strategy lift and action map, causal report transformation, a declared call/physical-time map, and a compatible audit request and task readout. Suppose the total variation distance of the complete stopped and scored laws is at most $\varepsilon$, uniformly over the target policies under comparison. An $M$-label target observer then gives a source observer with at most $RM$ labels and risk at most its target risk plus $L\varepsilon+e_{
m out}$, where $e_{
m out}$ is a separately certified expected output-loss perturbation. Serial compatible simulators multiply their label budgets, compose their call maps and add their complete-law errors. A reverse risk lower requires a reverse simulation with its own resource map.
\end{proposition}
\begin{proof}
Store the pair of simulator and observer labels, route legal actions through the strategy lift, generate each target report without future information, and apply the target update. Independence from the unknown physical parameter, if any is included in the experiment family, ensures that this is one program. The compatible request fixes the same terminal task; it is not replaced by a model-specific baseline. Integrating the bounded loss against the complete-law difference gives $L\varepsilon$, with total variation defined as the supremum of event differences. Add the output perturbation. Causal composition and the triangle inequality give the serial claim. A one-way inequality has no converse by itself.
\end{proof}

\begin{proposition}[Stationary perturbation and finite-prefix separation]\label{prop:perturbation}
Let $P$ be a fixed observer's positive ensemble map satisfying \eqref{eq:ensemble-minor} with $\gamma_*=\alpha\eta_*$. Let $\widetilde P$ be another positive mass-preserving finite-dimensional ensemble map and suppose
\[
 \sup_{Z\ge0,\ e_{\rm tot}Z=1}\norm{(P-\widetilde P)Z}_{\oplus B}\le\tau.
\]
Every invariant normalized ensemble $\widetilde z$ of $\widetilde P$ satisfies
\begin{equation}\label{eq:perturbation}
 \norm{z^*-\widetilde z}_{\oplus B}\le\frac\tau{\gamma_*}.
\end{equation}
If the perturbed program retains an independent erasure and the same raw minorization, its invariant ensemble is unique and attracts every initial ensemble. Fixed-readout stationary risk differs by at most $L\tau/\gamma_*+e_{
m out}$. By contrast a uniform per-step complete-law discrepancy $\tau$ gives only a finite-prefix bound $n\tau$ without an additional stability argument; it does not give infinite-transcript total variation convergence.
\end{proposition}
\begin{proof}
Existence of an invariant normalized ensemble follows from the finite-dimensional compact convex base and continuity. Subtract $z^*=Pz^*$ from $\widetilde z=\widetilde P\widetilde z$ and use the contraction of $P$ on equal-mass differences:
\[
 \norm{z^*-\widetilde z}\le(1-\gamma_*)\norm{z^*-\widetilde z}+\tau.
\]
This proves \eqref{eq:perturbation}. The perturbed raw erasure minorization proves its own attraction as before. The loss functional is bounded affine, so the risk bound follows. The finite-prefix statement follows by sequential coupling or telescoping kernels. It has a factor equal to the actual maximum number of compared calls; no limiting statement is inferred from it.
\end{proof}

\begin{corollary}[Attained calibration and deficiency in one protected task]\label{cor:joint}
There is a physical extension of every experiment in Theorem~\ref{thm:main} whose single summed-square task has an unobservable calibration alternative $\pm\delta$ and an attained early-bit deficiency $\varepsilon$, and whose worst-case calibration risk satisfies, for all sufficiently large $M$,
\begin{equation}\label{eq:joint}
 cM^{-2/d}-C/n+\delta^2+2\varepsilon
 \le\calR^{\rm ext}_{\rm hom}(n,M)-J_*(n)
 \le C\left(M^{-2/d}+\frac{M^{2/d}}n\right)+\delta^2+2\varepsilon.
\end{equation}
The maximum physical-call count is $N=n+3$. In the stationary limit, and also in the matching entrance region, this is a matching resource--resolution--risk law. Numerical and simulator errors are separately added under Propositions~\ref{prop:morphism}--\ref{prop:perturbation}, not converted automatically into lower bounds.
\end{corollary}
\begin{proof}
After the protected random request, freeze all core acquisition actions. Prepare an independent fair bit $B\in\{-1,1\}$ and send it through a report that reveals $B$ with probability $1-a$ and gives a marked erasure otherwise. This is one paid raw call. Require commitment before the final joint audit, which reveals both the core audit and $B$. On the binary parameter family of prescribed bit preparations, the deficiency of the erased report relative to the exact bit, with virtual report completion before the audit, is exactly $a/2$. Randomly fill an erased report for the upper; for the lower the sum of the two binary testing errors on an erased report is at least $a$. Hence put $\varepsilon=a/2$. The fair-bit Bayes squared risk is exactly $a=2\varepsilon$, using readouts $-1,0,1$. The audit cannot be consulted early by a simulator.

The calibration coordinate is an independent physical quantity fixed at either $+\delta$ or $-\delta$, invisible in all preceding reports. Its minimax squared risk is $\delta^2$: read zero for the upper, and average the two alternatives for the lower. It cannot be learned during the core. The extension's loss is the sum of the original audit loss and these two squared losses, so orthogonality/independence makes the lower additive. Conditional on the original observer label, the independent bit report cannot improve core prediction; arbitrary extra randomization does not alter the core lower.

For the upper use $m$ running core labels and, after the request, three readout copies for each of them, corresponding to the bit posterior $-1,0,1$. At most $4m+2$ labels implement all running, terminal and administrative roles. Take $m$ a fixed fraction of $M$. This is a sufficient construction, not a minimal-state theorem for the enlarged report alphabet. It preserves the core risk order, and the extra coordinate readouts give the exact additive terms. Preparing the core, at most $n$ core reports, the bit call and the single joint audit gives $N=n+3$. The bit is prepared after controlled actions stop; preparing it earlier would allow later actions to store it in the physical system and would invalidate this lower. The actual calibration alternative and erased-bit channel, rather than permitted error ceilings, are what justify the two added lower terms.
\end{proof}
